\documentclass{article}
\usepackage[spanish]{babel}
\usepackage{amsmath}

\title{Proyecto Final \\ Inteligencia Artificial para Agroindustria}
\author{Tu Nombre}
\date{\today}

\begin{document}

\maketitle

\section{Introducción}

Este proyecto presenta el desarrollo de dos modelos de clasificación utilizando técnicas de inteligencia artificial: uno aplicado a datos tabulares y otro a imágenes.

\section{Componente Tabular}

\subsection{Descripción del dataset}

El dataset utilizado contiene 4839 observaciones con 5 variables de entrada numéricas (V1–V5) y una variable objetivo denominada \textit{Class}. Este problema corresponde a una tarea de clasificación supervisada binaria.

La variable objetivo presenta dos clases, las cuales fueron codificadas como 0 y 1. Se observó un desbalance significativo entre clases, siendo una de ellas considerablemente mayoritaria.

\subsection{Preprocesamiento}

El preprocesamiento de los datos incluyó las siguientes etapas:

\begin{itemize}
    \item Conversión de la variable objetivo desde formato \texttt{bytes} a \texttt{string}.
    \item Codificación de la variable objetivo utilizando \textit{LabelEncoder}.
    \item Separación de variables de entrada (X) y variable objetivo (y).
    \item División del dataset en entrenamiento (80\%) y prueba (20\%) utilizando muestreo estratificado.
    \item Escalado de las variables mediante \textit{StandardScaler}, ajustado únicamente sobre el conjunto de entrenamiento.
    \item Conversión de los datos a tensores de PyTorch.
\end{itemize}

\subsection{Modelo}

Se implementó un modelo de red neuronal tipo perceptrón multicapa (MLP) con la siguiente arquitectura:

\begin{itemize}
    \item Capa de entrada: 5 variables
    \item Capa oculta 1: 16 neuronas con activación ReLU
    \item Capa oculta 2: 8 neuronas con activación ReLU
    \item Capa de salida: 2 neuronas (una por clase)
\end{itemize}

Matemáticamente, el modelo realiza una transformación lineal seguida de una función de activación no lineal:

\[
z = Wx + b
\]

y posteriormente genera logits que son utilizados por la función de pérdida.

No se aplica \textit{Softmax} en la salida, ya que la función \textit{CrossEntropyLoss} lo incorpora internamente.

\subsection{Entrenamiento}

El modelo fue entrenado durante 20 épocas utilizando el optimizador Adam con una tasa de aprendizaje de 0.001.

La función de pérdida utilizada fue \textit{CrossEntropyLoss}, adecuada para problemas de clasificación multiclase.

Durante el entrenamiento se observó una disminución progresiva de la función de pérdida, indicando que el modelo estaba aprendiendo la relación entre las variables de entrada y la variable objetivo.

\subsection{Resultados}

El modelo alcanzó una accuracy del 94.63\% en el conjunto de prueba.

La matriz de confusión obtenida fue:

\[
\begin{bmatrix}
916 & 0 \\
52 & 0
\end{bmatrix}
\]

\subsection{Análisis}

A pesar de la alta accuracy obtenida, la matriz de confusión evidencia un problema importante: el modelo clasifica todas las observaciones como la clase mayoritaria, sin lograr identificar correctamente la clase minoritaria.

Esto se debe al desbalance presente en el dataset, lo cual hace que la accuracy no sea una métrica suficiente para evaluar el desempeño del modelo.

Se intentó corregir este problema mediante el uso de pesos en la función de pérdida, sin embargo, el modelo continuó mostrando el mismo comportamiento.

Este resultado resalta la importancia de considerar métricas adicionales y técnicas específicas para el manejo de desbalance de clases en problemas de clasificación.

\section{Componente de Imágenes}

\subsection{Descripción del dataset}

El dataset utilizado contiene 456 imágenes de hojas de papa clasificadas en tres categorías:

\begin{itemize}
    \item Potato\_\_\_Early\_blight
    \item Potato\_\_\_Late\_blight
    \item Potato\_\_\_healthy
\end{itemize}

Este problema corresponde a una tarea de clasificación multiclase basada en imágenes.

\subsection{Preprocesamiento}

Las imágenes fueron transformadas mediante:

\begin{itemize}
    \item Redimensionamiento a $128 \times 128$
    \item Conversión a tensores
    \item Normalización en el rango $[-1,1]$
\end{itemize}

El dataset se dividió en 80\% para entrenamiento y 20\% para prueba.

\subsection{Modelo}

Se implementó una red neuronal convolucional (CNN) compuesta por dos bloques convolucionales seguidos de capas densas.

Cada bloque convolucional incluye:

\begin{itemize}
    \item Convolución 2D
    \item Función de activación ReLU
    \item MaxPooling
\end{itemize}

La red finaliza con capas densas que producen los logits de salida para las tres clases.

\subsection{Entrenamiento}

El modelo fue entrenado durante 10 épocas utilizando el optimizador Adam y la función de pérdida CrossEntropyLoss.

Se observó una disminución significativa de la pérdida durante el entrenamiento, indicando que el modelo aprendió patrones relevantes en las imágenes.

\subsection{Resultados}

El modelo alcanzó una accuracy del 80.43\% en el conjunto de prueba.

\subsection{Análisis}

La disminución rápida de la función de pérdida sugiere que el modelo se ajusta fuertemente a los datos de entrenamiento. Dado el tamaño reducido del dataset, esto podría indicar un posible sobreajuste.

Sin embargo, la accuracy obtenida demuestra que el modelo es capaz de capturar características relevantes para diferenciar entre las distintas clases de hojas de papa.

\section{Conclusiones}

Se desarrollaron modelos funcionales que permiten resolver problemas de clasificación en contextos agroindustriales, evidenciando tanto fortalezas como limitaciones en los datos y modelos utilizados.

\end{document}